%======================================================================
\chapter{Methods}
%======================================================================

\section{Methods from Neural Networks}

Mime is implemented as two coupled neural pipelines: a \textbf{speech translation pipeline} (STT $\rightarrow$ MT $\rightarrow$ TTS) and an \textbf{audio-to-blendshape (ABS) pipeline} that drives a 3D ARKit avatar in real time.

\subsection{Speech Translation Pipeline}

\noindent\textbf{Speech-to-Text (STT):} The microphone stream is captured at 16 kHz mono and segmented online using an RMS-based voice activity strategy. Chunks are finalized after a silence window and transcribed with Groq's \texttt{whisper\allowbreak-large\allowbreak-v3\allowbreak-turbo}. This chunked design reduces unnecessary API calls and keeps inference focused on complete utterances.

\noindent\textbf{Machine Translation (MT):} Each English utterance is translated to conversational French using \texttt{llama\allowbreak-3.3\allowbreak-70b\allowbreak-versatile} through a constrained system prompt. A short rolling context buffer is included to improve coherence for fragments and pronoun resolution during live dialogue.

\noindent\textbf{Text-to-Speech (TTS):} French text is synthesized by Inworld TTS (default model \texttt{inworld\allowbreak-tts\allowbreak-1.5\allowbreak-mini}, French voice). Audio is streamed as LINEAR16 at 48 kHz so playback can begin before sentence completion, which lowers perceived latency in meetings.

\subsection{Audio-to-Blendshape Neural Model}

The visual pipeline predicts 52 ARKit blendshape coefficients from audio. The training notebook uses a strided 1D convolutional frontend for local phonetic feature extraction, followed by:

\begin{enumerate}
    \item a temporal depthwise mixing block,
    \item sinusoidal positional encoding,
    \item a Transformer encoder for long-range temporal context,
    \item a regression head with sigmoid outputs to constrain coefficients to $[0,1]$.
\end{enumerate}

Given predicted sequence $\hat{Y}$ and target sequence $Y$, the core objective is \textit{masked mean squared error} over valid (unpadded) frames:

$$
\mathcal{L}_{\text{mse}} = \frac{1}{N}\sum_{t \in \Omega} \lVert \hat{Y}_t - Y_t \rVert_2^2
$$

To encourage realistic temporal dynamics and reduce jitter, training adds a velocity consistency term:

$$
\mathcal{L}_{\text{vel}} = \frac{1}{N'}\sum_{t \in \Omega'} \lVert (\hat{Y}_t-\hat{Y}_{t-1}) - (Y_t-Y_{t-1}) \rVert_2^2
$$

with total loss:

$$
\mathcal{L} = \mathcal{L}_{\text{mse}} + \lambda\,\mathcal{L}_{\text{vel}}, \quad \lambda = 0.35
$$

Training uses AdamW, OneCycleLR scheduling, mixed precision on CUDA, gradient clipping, and early stopping.

\subsection{Real-Time Inference and Rendering}

At runtime, the client uses an optimized ABS inference path with a sliding audio window and queue-based frame updates. Predicted blendshape vectors are remapped from BEAT ordering to canonical ARKit ordering, then applied to morph targets in a rigged GLB avatar via an offscreen renderer.

A transport bridge publishes:

\begin{enumerate}
    \item rendered avatar frames to a virtual camera, and
    \item synthesized TTS audio to a virtual cable device,
\end{enumerate}

allowing direct integration with Zoom-compatible conferencing workflows.

\section{Dataset and Data Processing}

\subsection{Primary Dataset}

The project is trained on the \textbf{BEAT multimodal talking dataset} \cite{liu2022beat}, using locally downloaded paired audio and blendshape annotations. In the current experiments, the active split comes from the English subset (\texttt{beat\_english\_v0.2.1}) available under the repository data directory.

\subsection{Preprocessing Pipeline}

Each sample is formed from a side-by-side \texttt{.wav} and \texttt{.json} pair:

\begin{enumerate}
    \item Audio is loaded as mono and resampled to 16 kHz.
    \item Per-sample clip duration is capped (10 seconds in the current notebook setup) for stable training throughput.
    \item JSON frame weights are parsed and normalized to a fixed 52-dimensional target vector per frame (pad/truncate when needed).
    \item Targets are aligned to the effective audio duration using 60 FPS frame logic.
    \item Variable-length sequences are padded per batch with explicit length tensors for masked losses.
\end{enumerate}

The training pipeline currently uses an 80/20 train/validation split with a fixed random seed for reproducibility.

\section{Validation Strategy and Experiments}

\subsection{Model-Level Validation}

Validation is performed every epoch with sequence masking to avoid padding bias. The primary tracked metrics are:

\begin{enumerate}
    \item \textbf{Masked MSE}: main optimization-aligned error metric.
    \item \textbf{Masked MAE}: interpretable average absolute coefficient error.
    \item \textbf{Threshold Accuracy}: fraction of coefficients within an absolute error threshold (default 0.05).
\end{enumerate}

TensorBoard logging is used to track train/validation loss, accuracy, MAE, learning rate, and trial summaries. Best checkpoints are selected by validation loss with early stopping patience and minimum-delta logic.

\subsection{Hyperparameter Experiments}

The notebook supports structured trial sweeps over learning rate, hidden dimension, Transformer depth/heads, dropout, batch size, and epoch count. Each trial writes isolated logs and checkpoints, enabling direct comparison of convergence behavior and best-epoch performance.

\subsection{System-Level Validation}

Beyond offline metrics, the system is validated in \textit{end-to-end live runs}:

\begin{enumerate}
    \item microphone speech $\rightarrow$ transcription,
    \item English-to-French translation,
    \item streamed French synthesis,
    \item synchronized avatar mouth motion and Zoom bridge output.
\end{enumerate}

Operational logs (for example first audio chunk latency, per-epoch ETA, and module-level status) are used to assess practical real-time behavior and robustness under device reconnect or stream interruptions.
